\section{Introdução e objetivos}
\label{sec:objetivos}

Um neurônio artificial faz duas coisas: soma entradas multiplicadas por
pesos e passa o resultado por uma função não linear. O amplificador
operacional faz as duas. O somador inversor soma sinais com ganhos
definidos por razões de resistores, então cada peso da rede pode virar
um resistor. E quando a saída do amplificador encosta no trilho de
alimentação, ela satura. Essa saturação é a função não linear. A
pergunta do experimento: dá para treinar uma rede neural completa e
construir o circuito inteiro com AmpOps e resistores?

Para responder, treinamos uma rede de cinco estágios
($64\!-\!48\!-\!32\!-\!24\!-\!16\!-\!10$) que reconhece os dígitos do
MNIST reduzidos para $8\times 8$ pixels, ou seja, 64 tensões de entrada
entre \SI{0}{\volt} e \SI{1}{\volt}. O treinamento acontece em software,
mas já usando as regras do circuito: a ativação é a saturação real do
somador e cada peso precisa existir como resistor comercial da série
E96. Depois do treino, um script converte os pesos em netlist SPICE,
com 260 amplificadores e 5451 resistores de peso. Na saída não tem
cálculo extra nem processador: cada dígito tem um comparador, e a
resposta da rede é o único comparador que acende.

A validação foi feita no LTspice em duas rodadas da rede completa, uma
com um modelo ideal de AmpOp e outra com o macromodelo comercial do
LM741, usando trilhos calibrados para compensar a saturação do 741. Um
terceiro amplificador, o LT1810, moderno e rápido, foi caracterizado em
bancadas menores e serve de contraponto ao 741 ao longo do relatório:
satura a milivolts dos trilhos, decide 25 vezes mais rápido e cobra o
preço em compensação e equilíbrio de polarização.
Todo o fluxo roda em Docker com o ferramental claw-spice do
repositório, cada simulação em um contêiner separado.

Os objetivos do experimento são: treinar a rede sob as restrições do
circuito e passar de \SI{97}{\percent} de acurácia no teste do MNIST;
construir a rede só com somadores inversores, inversores de trilho e um
comparador por dígito, e conferir no LTspice que o circuito reproduz o
modelo; e medir o tempo de resposta do circuito, comparando com a
inferência em software e com uma estimativa de custo de um chip.
